\chapter{Gewichtete Rankings}
\label{chap:weighted-rankings}

\begin{myremark}[Algorithmusart in diesem Kapitel]
    Funktionsart: \textbf{Priorisierung}. \\
    Umsetzungsart: \textbf{regelbasiert und selbstlernend}.
\end{myremark}

\begin{myremark}
    Dieses Kapitel behandelt Algorithmen, die nicht primär Kategorien
    oder Zahlen vorhersagen, sondern Elemente in eine \textbf{Reihenfolge} bringen.
\end{myremark}

\section{Grundidee: Ordnen statt klassifizieren}

Bei einer Priorisierung wird nicht gefragt \enquote{Welche Klasse?} oder \enquote{Welcher Wert?},
sondern \enquote{Was soll zuerst kommen?}. Das Ergebnis ist eine \textbf{Rangliste}.

Typische Beispiele sind Suchmaschinen, Newsfeeds, Lernplattformen oder Support-Tickets.
Die Qualität eines Priorisierungsalgorithmus hängt davon ab, wie sinnvoll die Reihenfolge ist.

\section{Gewichtete Priorisierung}
Die gewichtete Priorisierung ist eine einfache Methode, um mehrere Kriterien zu einem Score zusammenzufassen. Dieser Algorithmus kommt zum Beispiel bei der Bewertung von Lernquellen (welches Buch, welcher Artikel, welches Video ist am besten geeignet?), bei der Auswahl von Support-Tickets (welches Problem soll zuerst gelöst werden?) oder bei der Priorisierung von Aufgaben (welche To-Do hat die höchste Dringlichkeit?) zum Einsatz.

\begin{mydefinition}[Gewichtete Priorisierung]
    Bei einer gewichteten Priorisierung bekommt jedes Kriterium ein \textbf{Gewicht}.
    Aus den Teilwerten $z_1, z_2, \ldots, z_n$ und Gewichten $w_1, w_2, \ldots, w_n$ entsteht ein Score:
    \[
        \text{Score}(x) = w_1 z_1 + w_2 z_2 + \cdots + w_n z_n = \sum_{i=1}^n w_i z_i
    \]
    Oft gilt $\sum_i w_i = 1$ und alle Teilwerte sind auf $[0,1]$ normiert.
\end{mydefinition}

\begin{mycaution}
    Wenn Merkmale unterschiedliche Skalen haben, verzerren grosse Zahlen die Rangliste.
    Darum müssen die Werte vor dem Zusammenfassen oft normalisiert werden.
    Für negative Kriterien wie \enquote{Aufwand} oder \enquote{Kosten} wird der Wert invertiert,
    z.\,B. durch $1-z$.
\end{mycaution}

\begin{myexample}
    Drei Quellen für ein Schulprojekt werden nach vier Kriterien bewertet:
    Relevanz, Aktualität, Vertrauenswürdigkeit und geringer Aufwand.
    Die Gewichte seien $0.40$, $0.25$, $0.25$ und $0.10$.
\end{myexample}

\begin{myexercise}[Gewichtete Scores berechnen]
    Drei Lernquellen wurden bereits auf $[0,1]$ normiert:

    \begin{center}
        \begin{tabular}{lcccc}
            \toprule
            Quelle & Relevanz & Aktualität & Vertrauen & Aufwand \\
            \midrule
            A      & 0.85     & 0.50       & 0.75      & 0.30    \\
            B      & 0.75     & 0.85       & 0.70      & 0.20    \\
            C      & 0.80     & 0.65       & 0.95      & 0.60    \\
            \bottomrule
        \end{tabular}
    \end{center}

    Berechnen Sie den Score mit
    \[
        \text{Score} = 0.40\cdot R + 0.25\cdot A + 0.25\cdot V + 0.10\cdot(1-\text{Aufwand})
    \]
    und ordnen Sie die drei Quellen.
    \begin{myanswer}
        A: $0.40\cdot0.85 + 0.25\cdot0.50 + 0.25\cdot0.75 + 0.10\cdot0.70 = 0.7225$ \\
        B: $0.40\cdot0.75 + 0.25\cdot0.85 + 0.25\cdot0.70 + 0.10\cdot0.80 = 0.7675$ \\
        C: $0.40\cdot0.80 + 0.25\cdot0.65 + 0.25\cdot0.95 + 0.10\cdot0.40 = 0.7600$ \\
        Rangfolge: B, C, A.
    \end{myanswer}
\end{myexercise}

\begin{myexercise}[Gewichte diskutieren]
    Was verändert sich, wenn man Aktualität stärker gewichtet als Vertrauenswürdigkeit?
    Nennen Sie ein Beispiel, in dem diese Entscheidung sinnvoll sein kann.
    \begin{myanswer}
        Wenn Aktualität wichtiger ist, könnte Quelle C vor B liegen, da C aktueller ist als B. Das könnte sinnvoll sein, wenn es um ein schnelllebiges Thema geht, z.\,B. aktuelle Nachrichten oder Trends.
    \end{myanswer}
\end{myexercise}